\documentclass[12pt, a4paper]{article}

\usepackage[]{graphicx}

% FONT
\usepackage{newtxtext}

\usepackage[margin=3cm]{geometry}

% interlinea
\usepackage{setspace}
\renewcommand{\baselinestretch}{1.5} 

\newcommand{\meskip}{\medskip \\}

\usepackage{enumitem}

\title{
    \includegraphics[width=.8\textwidth]{images/LOGO.png}\\
    \textbf{Executive Summary}
}
\date{}
\author{}

\begin{document}

\maketitle

\section*{Descrizione}
Benvenuto nel meraviglioso mondo del vino, dove la passione e l'eleganza si fondono in una bottiglia! La nostra startup, chiamata Vinovo, è nata con l'obiettivo di portare un'esperienza enologica eccezionale a tutti gli amanti del vino. \meskip
La nostra missione è quella di creare un legame tra i produttori di vino di alta qualità e i consumatori che cercano sapori raffinati.\meskip
Inoltre, la nostra piattaforma online offre ai nostri clienti la possibilità esplorare il nostro vasto assortimento di vini, leggere recensioni autentiche e ordinare comodamente da casa loro.
Offriamo anche un servizio di spedizione veloce e sicuro per garantire che ogni bottiglia raggiunga il cliente in condizioni ottimali.\meskip
Vinovo si impegna a promuovere la sostenibilità nel settore vinicolo.
Collaboriamo con produttori che adottano pratiche ecologiche nella produzione dei loro vini.
Vogliamo essere un punto di riferimento per i consumatori che si preoccupano dell'ambiente.\meskip
In sintesi, Vinovo è molto più di una semplice startup per la vendita di vini.
Siamo un'esperienza enologica completa, dove la qualità e la sostenibilità si incontrano per offrire ai nostri clienti un viaggio indimenticabile nel mondo del vino.
Unisciti a noi e scopri l'arte e il piacere di un buon bicchiere di vino.

\section*{Mercato}
Il settore del vino è un mercato ampio e in continua crescita, con una domanda stabile e una cultura consolidata intorno al consumo di vino.
Ci sono molte opportunità nel settore per offrire selezioni curate, servizi di consulenza e un'esperienza d'acquisto coinvolgente.\meskip
Il settore del vino è altamente competitivo, con numerosi rivenditori e piattaforme online che offrono una vasta gamma di vini.
I concorrenti includono enoteche tradizionali, negozi di vini specializzati, grandi catene di distribuzione e piattaforme di e-commerce dedicate al vino.\meskip
La sfida principale è differenziarsi dai concorrenti offrendo un prodotto esclusivo e una selezione di vini unica, esperienza personalizzata, prezzi competitivi e una forte presenza online.\meskip
C'è una crescente tendenza verso l'acquisto di vino online, con sempre più persone che vogliono scoprire nuove etichette provenienti da diverse regioni vinicole; inoltre i consumatori sono sempre più interessati a vini biologici e sostenibili.\meskip
I potenziali clienti includono gli amanti del vino, ristoratori, bar e hotel, nonché persone che cercano regali di alta qualità e unici legati al vino.
\newpage

\section*{Modello di business}
\textbf{Canali di distribuzione}:
\begin{itemize}[itemsep=-5pt, topsep=0pt]
    \item negozio online
    \item collaborazioni con ristoranti, hotel bar ed enoteche
    \item partecipare a eventi e fiere del settore
\end{itemize}\medskip
\textbf{Relazioni con i produttori}: si sviluppano relazioni con i produttori di vino artigianali, stabilendo accordi diretti per l'acquisto delle loro produzioni. Questo permette di garantire la qualità dei vini.\meskip
\textbf{Marketing e promozione}: si utilizzano strategie di marketing mirate per promuovere l'offerta di vini come la creazione di contenuti sul proprio sito web, l'utilizzo dei social media per condividere informazioni sulle etichette e le storie dei produttori, partecipazione a eventi di degustazione.\meskip
\textbf{Fidelizzazione dei clienti}: per incentivare i clienti fedeli a continuare ad acquistare, si implementano sconti speciali, accesso anticipato a nuove selezioni o eventi esclusivi.\meskip
\textbf{Logistica}: si implementa un sistema logistico efficiente per garantire la consegna rapida e sicura dei vini ai clienti.\meskip
\textbf{Scalabilità}: si pianifica una crescita graduale, aumentando gradualmente l'offerta di vini e allargando la sua base di clienti. Potrebbe anche esplorare opportunità di espansione geografica o di offerta di servizi di abbonamento per creare entrate ricorrenti.
\newpage
\section*{Analisi del piano finanziario}
\begin{tabular}{|c|c|c|c|}
    \hline
                & 1° anno & 2° anno      & 3° anno        \\ \hline
    R.O.I. (\%) & PERDITA & 4\%          & 72\%           \\ \hline
    R.O.E. (\%) & PERDITA & 1\%          & 78\%           \\ \hline
    BEP         & PERDITA & 610.671,30 € & 1.164.468,30 € \\ \hline
\end{tabular}\bigskip
$$VAN = 132.884,05 \; €$$

\section*{Cosa ci rende diversi?}
In primo luogo, offriamo vini artigianali di nostra produzione, offrendo così un prodotto esclusivo e non reperibile altrove.\\
Inoltre, ci impegniamo a offrire una selezione di vini provenienti dalle rinomate regioni vinicole italiane, garantendo così prodotti di qualità.\meskip
Tuttavia, Vinovo non si limita solo a essere una piattaforma di vendita di vini pregiati.
Mettiamo a disposizione guide dettagliate, consigli sull'abbinamento cibo-vino e approfondimenti sulle diverse varietà di uve e regioni vinicole.

\end{document}
